\documentclass[11pt]{amsart}
\usepackage{amsmath,amssymb,amsthm}
\usepackage[margin=1.1in]{geometry}
\usepackage[colorlinks=true,allcolors=blue]{hyperref}
\usepackage[capitalise,nameinlink]{cleveref}

\newtheorem{theorem}{Theorem}[section]
\newtheorem{proposition}[theorem]{Proposition}
\newtheorem{lemma}[theorem]{Lemma}
\newtheorem{corollary}[theorem]{Corollary}
\theoremstyle{definition}
\newtheorem{definition}[theorem]{Definition}
\theoremstyle{remark}
\newtheorem{remark}[theorem]{Remark}

\newcommand{\E}{\mathbb{E}}
\newcommand{\PP}{\mathbb{P}}
\newcommand{\R}{\mathbb{R}}
\newcommand{\Rad}{R}

\title[Randomized smoothing at high confidence]{Randomized smoothing at high confidence: the Gaussian is optimal for every noise law}
\author{Yitzchak Shmalo}
\date{\today}

\begin{document}

\begin{abstract}
Randomized smoothing certifies a classifier by averaging a base classifier over a noise law $q$ on $\R^d$. We
consider certificates that hold for every base classifier and require only that the top class has probability
at least $1-\alpha$ under the noise. We prove that if $\E\|\eta\|_2^2\le d\sigma^2$ for $\eta\sim q$, then for
$0<\alpha<1/4$ and $d\ge 1012\,\alpha^{-2}\log(16/\alpha)$ no such certificate exceeds
$8\sigma\sqrt{\log(16/\alpha)}\,d^{1/p-1/2}$ in the $\ell_p$ norm, $2\le p\le\infty$. The Gaussian
$N(0,\sigma^2 I)$ certifies $\sigma\Phi^{-1}(1-\alpha)\,d^{1/p-1/2}$, so it is optimal up to an absolute
constant factor at every confidence level. Earlier bounds for general noise laws allow the certified radius at
fixed noise level to grow like $\alpha^{-1/2}$ as $\alpha\to0$; the correct rate is $(\log 1/\alpha)^{1/2}$.
The proof uses only the second moment of $q$: no product structure, symmetry, density or log-concavity.
\end{abstract}

\maketitle

\section{Introduction}

Let $f\colon\R^d\to\mathcal{Y}$ be a measurable classifier and let $q$ be a Borel probability measure on
$\R^d$. The smoothed classifier $g(x)=\arg\max_{c}\PP_{\eta\sim q}(f(x+\eta)=c)$ is the object of
randomized smoothing \cite{LAG19,LCWC19,CRK19}. For the Gaussian $q=N(0,\sigma^2I)$, Cohen, Rosenfeld and
Kolter \cite{CRK19} showed by the Neyman--Pearson lemma that if the top class $A$ has probability $p_A>1/2$
at $x$, then $g(x+\delta)=A$ whenever $\|\delta\|_2<\sigma\Phi^{-1}(p_A)$, and that this radius cannot be
improved using $p_A$ alone. Since the $\ell_\infty$ ball of radius $r$ lies in the $\ell_2$ ball of radius
$r\sqrt d$, the same noise certifies the $\ell_\infty$ radius $\sigma\Phi^{-1}(p_A)/\sqrt d$.

Whether another noise law does better for $\ell_p$, $p>2$, has been studied from several directions. Blum,
Dick, Manoj and Zhang \cite{BDMZ20} proved that a noise law which certifies $\ell_p$ radius $\varepsilon$ for
every base classifier with score gap $\delta$ must have $\E\eta_i^2=\Omega(d^{1-2/p}\varepsilon^2
(1-\delta)/\delta^2)$ for most coordinates $i$. Kumar, Levine, Goldstein and Feizi \cite{KLGF20} showed that
for a large class of product noise laws the certified $\ell_p$ radius decreases as $O(d^{1/p-1/2})$. Yang et
al.\ \cite{YDHSRL20} obtained a bound of the same order in $d$ for a broad class of smoothing schemes that use
only the probability of the top class, including schemes whose noise depends on the input, and Wu,
Bojchevski, Kuvshinov and G\"unnemann \cite{WBKG21} gave upper bounds of order $d^{-1/2}$ for the $\ell_2$
radius under continuous symmetric noise.

Read at a fixed noise level $\sigma$ per coordinate and confidence $p_A=1-\alpha$, so that the score gap is
$1-2\alpha$, the bound of \cite{BDMZ20} gives $\varepsilon=O(\sigma d^{1/p-1/2}\alpha^{-1/2})$, while the
Gaussian certifies $\sigma\Phi^{-1}(1-\alpha)d^{1/p-1/2}\sim\sigma\sqrt{2\log(1/\alpha)}\,d^{1/p-1/2}$. The
gap between $\alpha^{-1/2}$ and $(\log1/\alpha)^{1/2}$ is the question of whether a non-Gaussian law can buy
more radius at high confidence. The theorem below shows it cannot.

\begin{definition}\label{def:radius}
Let $\|\cdot\|$ be a norm on $\R^d$ and $0<\alpha<1/2$. The \emph{certified radius} $\Rad_q(\alpha,\|\cdot\|)$
is the supremum of $\varepsilon\ge0$ with the following property: for every measurable
$f\colon\R^d\to\mathcal{Y}$, every $x\in\R^d$ and every class $A$ with $\PP_{\eta\sim q}(f(x+\eta)=A)\ge1-\alpha$,
one has $\PP_{\eta\sim q}(f(x+\delta+\eta)=A)>1/2$ for all $\|\delta\|\le\varepsilon$.
\end{definition}

The condition $\PP(f(x+\delta+\eta)=A)>1/2$ forces $g(x+\delta)=A$, so $\Rad_q$ is the largest radius that can
be certified from the confidence level alone, for every base classifier. By the Cohen--Rosenfeld--Kolter bound,
$\Rad_{N(0,\sigma^2I)}(\alpha,\|\cdot\|_2)=\sigma\Phi^{-1}(1-\alpha)$.

\begin{theorem}\label{thm:main}
Let $q$ be a Borel probability measure on $\R^d$ with $\E_{\eta\sim q}\|\eta\|_2^2\le d\sigma^2$, where
$\sigma>0$. Let $0<\alpha<1/4$ and
\[
d\ \ge\ \frac{1012\,\log(16/\alpha)}{\alpha^{2}} .
\]
Then for every $2\le p\le\infty$,
\[
\Rad_q(\alpha,\|\cdot\|_p)\ \le\ 8\,\sigma\sqrt{\log(16/\alpha)}\;d^{1/p-1/2}.
\]
\end{theorem}

\begin{corollary}\label{cor:gauss}
Under the assumptions of \cref{thm:main}, for every $2\le p\le\infty$,
\[
\Rad_q(\alpha,\|\cdot\|_p)\ \le\ 25\,\Rad_{N(0,\sigma^2I)}(\alpha,\|\cdot\|_p),
\]
and the ratio of the two bounds tends to at most $8/\sqrt2<5.7$ as $\alpha\to0$.
\end{corollary}

The constants are not optimized. Only the second moment of $q$ enters; the noise may be any law, with atoms,
heavy tails or strong dependence between coordinates. The proof is an explicit construction: a base
classifier that has probability at least $1-\alpha$ of the top class at $x$ and probability at most $1/2$ at a
point $x+v$ with $\|v\|_\infty=\varepsilon$. The shift $v$ is a multiple of a sign vector, which is why the
$\ell_p$ bounds follow from the $\ell_\infty$ one.

\section{A criterion for failure}

\begin{lemma}\label{lem:criterion}
Let $\|\cdot\|$ be a norm on $\R^d$. Suppose there are a Borel set $B_0\subset\R^d$ and a vector $v$ with
$q(B_0)\ge1/2$ and $q(v+B_0)\le\alpha$. Then $\Rad_q(\alpha,\|\cdot\|)\le\|v\|$.
\end{lemma}

\begin{proof}
Fix $x$ and two classes $A\ne B$. Let $f(y)=B$ if $y\in x+v+B_0$ and $f(y)=A$ otherwise. Then
$\PP(f(x+\eta)=B)=q(v+B_0)\le\alpha$, so $\PP(f(x+\eta)=A)\ge1-\alpha$. At the point $x+v$,
$\PP(f(x+v+\eta)=B)=q(B_0)\ge1/2$, so $\PP(f(x+v+\eta)=A)\le1/2$. Hence no $\varepsilon\ge\|v\|$ has the
property of \cref{def:radius}.
\end{proof}

\section{Proof of the theorem}

Throughout, $\eta\sim q$, $B_r=\{y:\|y\|_2\le r\}$, and $\varepsilon>0$ is a number with $\varepsilon\le\sigma$ to
be fixed at the end.

\begin{lemma}[a thin shell]\label{lem:shell}
There is $\rho\in[2\sigma\sqrt d,\,3\sigma\sqrt d-\varepsilon\sqrt d\,]$ with $q(B_\rho)\ge3/4$ and
\[
q\bigl(\rho<\|\eta\|_2\le\rho+\varepsilon\sqrt d\bigr)\ \le\ \frac{2\varepsilon}{\sigma}.
\]
\end{lemma}

\begin{proof}
By Markov's inequality $q(\|\eta\|_2>2\sigma\sqrt d)\le\E\|\eta\|_2^2/(4d\sigma^2)\le1/4$, so $q(B_\rho)\ge3/4$
for every $\rho\ge2\sigma\sqrt d$. The interval $[2\sigma\sqrt d,3\sigma\sqrt d]$ contains
$N=\lfloor\sigma/\varepsilon\rfloor\ge1$ disjoint half-open intervals $(\rho_k,\rho_k+\varepsilon\sqrt d]$,
$\rho_k=2\sigma\sqrt d+k\varepsilon\sqrt d$, $0\le k<N$. The corresponding shells are disjoint, so one of them
has mass at most $1/N\le2\varepsilon/\sigma$, using $\lfloor t\rfloor\ge t/2$ for $t\ge1$.
\end{proof}

Fix $\rho$ as in \cref{lem:shell} and let $\mu$ be the law of $\eta$ conditioned on $\eta\in B_\rho$. For a sign
vector $s\in\{-1,1\}^d$ write $X_s(y)=\langle y,s\rangle/\sqrt d$. Put
\[
T=\sigma\sqrt{18\log24},\qquad T'=\sigma\sqrt{18\log(16/\alpha)} .
\]

\begin{lemma}[a direction]\label{lem:direction}
There is $s\in\{-1,1\}^d$ with $\mu(X_s>T)\le1/3$ and $\mu(X_s<-T')\le\alpha/2$.
\end{lemma}

\begin{proof}
Let $s$ be uniform on $\{-1,1\}^d$, independent of $\eta\sim\mu$. For fixed $y\in B_\rho$, $\langle y,s\rangle$ is a
sum of independent bounded terms $s_iy_i$, and Hoeffding's inequality gives, for $u>0$,
\[
\PP_s\bigl(X_s(y)>u\bigr)\le\exp\Bigl(-\frac{u^2d}{2\|y\|_2^2}\Bigr)\le\exp\Bigl(-\frac{u^2}{18\sigma^2}\Bigr),
\]
since $\|y\|_2\le\rho\le3\sigma\sqrt d$; the same bound holds for $\PP_s(X_s(y)<-u)$. Integrating over $y$,
$\E_s\,\mu(X_s>T)\le1/24$ and $\E_s\,\mu(X_s<-T')\le\alpha/16$. By Markov's inequality in $s$,
$\PP_s(\mu(X_s>T)>1/3)\le1/8$ and $\PP_s(\mu(X_s<-T')>\alpha/2)\le1/8$, so some $s$ satisfies both bounds.
\end{proof}

\begin{proof}[Proof of \cref{thm:main}]
Take $s$ from \cref{lem:direction} and set
\[
B_0=\{y\in B_\rho:\ X_s(y)\le T\},\qquad \varepsilon=\frac{T+T'}{\sqrt d},\qquad v=-\varepsilon s .
\]
First, $q(B_0)=q(B_\rho)\,\mu(X_s\le T)\ge\frac34\cdot\frac23=\frac12$.

Next let $y=v+z$ with $z\in B_0$. Then $\|y\|_2\le\rho+\|v\|_2=\rho+\varepsilon\sqrt d$ and
$X_s(y)=X_s(z)-\varepsilon\sqrt d\le T-(T+T')=-T'$. So $y$ lies either in the shell
$\{\rho<\|y\|_2\le\rho+\varepsilon\sqrt d\}$ or in $B_\rho\cap\{X_s\le-T'\}$, and
\[
q(v+B_0)\ \le\ \frac{2\varepsilon}{\sigma}+q(B_\rho)\,\mu(X_s\le-T')\ \le\ \frac{2\varepsilon}{\sigma}+\frac{\alpha}{2},
\]
where $\mu(X_s\le-T')\le\alpha/2$ holds by the proof of \cref{lem:direction}, the Hoeffding bound being the same
for the closed event.

It remains to check $2\varepsilon/\sigma\le\alpha/2$ and $\varepsilon\le\sigma$. For $\alpha<1/4$ we have
$\log(16/\alpha)>\log64$, hence
\[
\frac{T+T'}{\sigma}=\sqrt{18}\Bigl(\sqrt{\log24}+\sqrt{\log(16/\alpha)}\Bigr)
\le\sqrt{18}\Bigl(\sqrt{\tfrac{\log24}{\log64}}+1\Bigr)\sqrt{\log(16/\alpha)}\le 7.952\sqrt{\log(16/\alpha)} .
\]
The condition $\varepsilon\le\alpha\sigma/4$ reads $\sqrt d\ge4(T+T')/(\alpha\sigma)$, which holds when
$d\ge16\cdot7.952^2\,\alpha^{-2}\log(16/\alpha)$, and $16\cdot7.952^2<1012$. Then also $\varepsilon\le\sigma$.
So $q(v+B_0)\le\alpha$, and \cref{lem:criterion} gives, since $v$ is $\varepsilon$ times a sign vector,
\[
\Rad_q(\alpha,\|\cdot\|_p)\le\|v\|_p=\varepsilon\,d^{1/p}=\frac{T+T'}{\sigma}\,\sigma\,d^{1/p-1/2}
\le8\,\sigma\sqrt{\log(16/\alpha)}\,d^{1/p-1/2}.\qedhere
\]
\end{proof}

\begin{proof}[Proof of \cref{cor:gauss}]
For $p\ge2$, $\|\delta\|_2\le d^{1/2-1/p}\|\delta\|_p$, so the Gaussian certifies the $\ell_p$ radius
$\sigma\Phi^{-1}(1-\alpha)\,d^{1/p-1/2}$. The function $\alpha\mapsto8\sqrt{\log(16/\alpha)}/\Phi^{-1}(1-\alpha)$ is
at most $24.19$ on $(0,1/4)$, with the largest values near $\alpha=1/4$, and tends to $8/\sqrt2$ as $\alpha\to0$
because $\Phi^{-1}(1-\alpha)\sim\sqrt{2\log(1/\alpha)}$.
\end{proof}

\section{Remarks}

\begin{remark}
The Gaussian value $\sigma\Phi^{-1}(1-\alpha)/\sqrt d$ in $\ell_\infty$ is exact for $N(0,\sigma^2I)$: the shift
$v=\varepsilon(1,\dots,1)$ and the half-space classifier $f(y)=B$ on $\{\langle y-x,(1,\dots,1)\rangle\ge
\sigma\sqrt d\,\Phi^{-1}(1-\alpha)\}$ attain it. \Cref{thm:main} says that no other law with the same average
per-coordinate second moment improves this by more than a constant factor, at any confidence level.
\end{remark}

\begin{remark}
The noise level is measured by the average second moment $d^{-1}\E\|\eta\|_2^2$, as in \cite{KLGF20}; the bound
of \cite{BDMZ20} is stated for most individual coordinates.
\end{remark}

\begin{remark}
The condition $d\ge C\alpha^{-2}\log(1/\alpha)$ comes from the shell step, which needs a shell of relative
mass $O(\alpha)$, and is not optimized. Three questions remain: whether the Gaussian is exactly optimal among
laws with a given second moment; whether the condition on $d$ can be removed; and what happens for
input-dependent smoothing, where the noise law varies with $x$ \cite{SKH21}. The argument above uses a single
law $q$ at $x$ and at $x+v$ and does not apply directly to the last case.
\end{remark}

\section*{Acknowledgments and funding}
This research was supported by the European Research Council (ERC) under the European Union's Horizon Europe
research and innovation programme (grant agreement No.~101041711), by the Simons Foundation as part of the
Collaboration on the Mathematical and Scientific Foundations of Deep Learning, by Heights Labs, by Convex Nexus
Capital, by the Israel Science Foundation (grant number 2258/19), and by the Israel Science Foundation (ISF Grant
4101/25).

\section*{Declaration of competing interest}
The author holds equity in Heights Labs and in Convex Nexus Capital and has a management role at Convex Nexus
Capital. These entities are included in the funding declaration above.

\section*{Declaration of generative AI and AI-assisted technologies}
The paper comes out of Math Brain, an ongoing project on the use of AI systems for research mathematics. In
November 2026 Claude models (Anthropic), working in Claude Code, proved the theorem, checked the construction
numerically, and wrote the paper. The author directed the work and is responsible for the content of the paper
and for decisions concerning its dissemination.

\begin{thebibliography}{99}

\bibitem{BDMZ20}
A.~Blum, T.~Dick, N.~Manoj and H.~Zhang, Random smoothing might be unable to certify $\ell_\infty$ robustness
for high-dimensional images, \emph{J. Mach. Learn. Res.} 21 (2020), no.~211, 1--21. arXiv:2002.03517.

\bibitem{CRK19}
J.~Cohen, E.~Rosenfeld and Z.~Kolter, Certified adversarial robustness via randomized smoothing, in
\emph{Proceedings of the 36th International Conference on Machine Learning}, PMLR 97 (2019), 1310--1320.

\bibitem{KLGF20}
A.~Kumar, A.~Levine, T.~Goldstein and S.~Feizi, Curse of dimensionality on randomized smoothing for certifiable
robustness, in \emph{Proceedings of the 37th International Conference on Machine Learning}, PMLR 119 (2020),
5458--5467. arXiv:2002.03239.

\bibitem{LAG19}
M.~Lecuyer, V.~Atlidakis, R.~Geambasu, D.~Hsu and S.~Jana, Certified robustness to adversarial examples with
differential privacy, in \emph{2019 IEEE Symposium on Security and Privacy}, 656--672.

\bibitem{LCWC19}
B.~Li, C.~Chen, W.~Wang and L.~Carin, Certified adversarial robustness with additive noise, in \emph{Advances in
Neural Information Processing Systems} 32 (2019).

\bibitem{SKH21}
P.~S\'uken\'ik, A.~Kuvshinov and S.~G\"unnemann, Intriguing properties of input-dependent randomized smoothing,
in \emph{Proceedings of the 39th International Conference on Machine Learning}, PMLR 162 (2022). arXiv:2110.05365.

\bibitem{WBKG21}
Y.~Wu, A.~Bojchevski, A.~Kuvshinov and S.~G\"unnemann, Completing the picture: randomized smoothing suffers from
the curse of dimensionality for a large family of distributions, in \emph{Proceedings of the 24th International
Conference on Artificial Intelligence and Statistics}, PMLR 130 (2021).

\bibitem{YDHSRL20}
G.~Yang, T.~Duan, J.~E.~Hu, H.~Salman, I.~Razenshteyn and J.~Li, Randomized smoothing of all shapes and sizes, in
\emph{Proceedings of the 37th International Conference on Machine Learning}, PMLR 119 (2020), 10693--10705.

\end{thebibliography}

\end{document}
